\begin{abstract}
课程平台单次训练只有 \qtyrange{25}{60}{\minute}，在这一预算下，王者荣耀 1v1 的平台近端策略优化（PPO）基线未能推塔。为此，把鲁班七号的战术知识写成采样前的合法动作掩码，以先验知识缩小策略的假设空间；掩码在训练与推理中共用并随样本保存，使新旧策略在同一动作支持集上计算重要性比率。方法另含混合塑形奖励与冷启动偏置。在腾讯开悟 V63.0.9 平台上开展三组同代码、同超参数的 \qty{30}{\minute} 冷启动训练，对内置规则对手的终点监控胜率为：完整方法 \num{1.00}，仅用稀疏奖励 \num{0.80}，去掉掩码 \num{0}；官方基线训练 \qty{25}{\minute} 后为 \num{0}。仅用稀疏奖励时胜率在训练开始即为 \num{0.67}，推塔能力主要来自掩码；混合奖励使己方外塔余血由 \num{3526} 升至 \num{8107}。每组配置只运行一次，结论的稳健性仍需多随机种子实验检验。
\end{abstract}

\begin{enabstract}
Under the 25--60 minute training budget of a course platform, the platform's proximal policy optimization (PPO) baseline for Honor of Kings 1v1 did not learn to push towers. Tactical knowledge for the hero Luban No.~7 is encoded as a legal-action mask applied before sampling, which uses prior knowledge to shrink the hypothesis space of policies. The mask is shared by inference and training and stored with each sample, so old and new policies compute importance ratios over the same action support. A shaped reward and a cold-start bias complete the method. In three 30-minute cold-start runs on the Tencent Kaiwu V63.0.9 platform with identical code and hyperparameters, final monitored win rates against the built-in rule-based opponent were 1.00 for the full method, 0.80 with sparse reward only, and 0 without the mask; the official baseline reached 0 after 25 minutes. With sparse reward only, the win rate was already 0.67 at the start of training, so tower pushing came mainly from the mask; the shaped reward raised the remaining hit points of the agent's own tower from 3526 to 8107. Each configuration was run once, and multi-seed experiments are needed to confirm robustness.
\end{enabstract}
